% Overlap-add: x em blocos de L, cada um convoluido na placa com h (M amostras).
\begin{tikzpicture}[
    blk/.style={draw, minimum height=6mm, inner sep=0pt, font=\scriptsize},
  ]
  \def\L{30mm}\def\T{10mm} % comprimento do bloco e da cauda (M-1) no desenho

  % Sinal de entrada
  \node[anchor=east, font=\small] at (-2mm,0) {\(x[n]\)};
  \foreach \i/\c in {0/opConv, 1/opFft, 2/opIir} {
    \node[blk, fill=\c!15, minimum width=\L, anchor=west] at (\i*\L,0) {bloco \the\numexpr\i+1\relax: \(L\) amostras};
  }
  \node[font=\small] at (3*\L+6mm,0) {\(\cdots\)};

  % Convolucoes na placa (deslocadas)
  \foreach \i/\c in {0/opConv, 1/opFft, 2/opIir} {
    \pgfmathsetmacro{\y}{-12-9*\i}
    \node[blk, fill=\c!15, minimum width=\L, anchor=west] (c\i) at (\i*\L,\y mm) {\(x_{\the\numexpr\i+1\relax} * h\)};
    \node[blk, fill=\c!35, minimum width=\T, anchor=west] at (c\i.east) {cauda};
    \node[font=\scriptsize, anchor=east, text=black!60] at (-2mm,\y mm) {na placa};
  }

  % Soma
  \draw[thick] (0,-42mm) -- (3*\L+\T,-42mm);
  \node[anchor=east, font=\small] at (-2mm,-48mm) {\(y[n]\)};
  \node[blk, fill=black!5, minimum width=3*\L+\T, anchor=west] at (0,-48mm)
    {soma: onde duas caudas se sobrepõem, as amostras se somam};

  % Marcas de sobreposicao
  \foreach \i in {1,2}
    \draw[densely dashed, black!50] (\i*\L,-8mm) -- (\i*\L,-45mm);
  \draw[decorate, decoration={brace, amplitude=3pt}]
    (\L,-37mm) -- node[below=3pt, font=\scriptsize]{\(M-1\)} (\L+\T,-37mm);
\end{tikzpicture}
